\chapter{DSL Design}
\label{ch:dsl-design}

The source language describes polynomial hash functions close to their mathematical presentation. It is intentionally small: expressions, functions, conditionals, buffer accesses, field arithmetic, and a few reductions. The examples are motivated by concrete designs, but the language is not hard-wired to them.

\section{Design Principles}
\label{sec:dsl-principles}

DSL functions are expression-oriented. A programmer writes the computed value, not a sequence of statements updating local variables. This mirrors polynomial hash definitions as sums, products, recurrences, and polynomial evaluations.

The language has no user-visible mutable variables. Mutable state appears only after lowering, for example as an accumulator in the IR. This simplifies type checking, code generation, and reference-test generation.

The DSL still contains buffers and index expressions, because generated hash functions read bytes from memory. It also distinguishes hash functions from helper functions: hash functions have a stable C API, while helper functions may return field elements or indices internally.

\section{Types}
\label{sec:dsl-types}

The language has three user-facing types.

\begin{description}
    \item[Index.] Index values are scalar integer values used for loop bounds, offsets, and length computations.
    \item[Field element.] Field elements are values in the selected prime field \(\mathbb{F}_{2^\pi-\theta}\). The source program does not fix \(\pi\) and \(\theta\); these are compiler parameters.
    \item[Buffer.] Buffers represent byte arrays passed at the C boundary. Loading from a buffer at field-element index \(i\) decodes the corresponding byte chunk into a field element.
\end{description}

Field elements are abstract in the source language. Limb count, limb width, and SIMD layout are introduced after the target field and platform are known.

\section{Core Grammar}
\label{sec:dsl-grammar}

\Cref{lst:dsl-grammar} gives a compact grammar for the central language. It omits lexical details, but shows the relevant structure: a module contains functions, every function body is one expression, and reductions are expressions rather than statements.

\begin{lstlisting}[style=thesisBNF,caption={Core grammar of the DSL.},label={lst:dsl-grammar}]
module      ::= function+
function    ::= hashfun | helperfun

hashfun     ::= "hash" ident "(" ident "buffer" ","
                              ident "buffer" ","
                              ident "index" ")" "{" expr "}"

helperfun   ::= "func" ident "(" params? ")" rettype "{" expr "}"
params      ::= paramgroup ("," paramgroup)*
paramgroup  ::= ident ("," ident)* type
rettype     ::= "fieldelement" | "index"
type        ::= "buffer" | "fieldelement" | "index"

expr        ::= comparison
comparison  ::= addsub (("=="|"!="|"<"|"<="|">"|">=") addsub)?
addsub      ::= muldiv (("+"|"-") muldiv)*
muldiv      ::= unary (("*"|"/"|"%") unary)*
unary       ::= "-" unary | power
power       ::= primary ("**" power)?

primary     ::= literal ("as" type)?
              | ident
              | ident "[" expr "]"
              | ident "(" (expr ("," expr)*)? ")"
              | "if" expr "{" expr "}" "else" "{" expr "}"
              | "nctif" expr "{" expr "}" ("else" "{" expr "}")?
              | reduction
              | "(" expr ")"

reduction   ::= sum | horner | foldl
sum         ::= "sum" "[" ident "," expr ":" expr ":" expr "]"
                "{" expr "}"
horner      ::= "horner" "[" ident "," expr ":" expr ":" expr
                          "," expr "]" "{" expr "}"
foldl       ::= "foldl" "[" ident "," expr ":" expr ":" expr
                         "," ident "=" expr "]" "{" expr "}"
\end{lstlisting}

\section{Hash Functions and Helper Functions}
\label{sec:hash-helper-functions}

A hash function is the external entry point. In generated C, every hash function has the same shape:
\begin{lstlisting}[style=thesisC]
void h(uint8_t *output, uint8_t *key,
       uint8_t *message, size_t len);
\end{lstlisting}
At this boundary, \(len\) is a byte length. The compiler converts it to a number of field elements.

Inside the DSL, the corresponding length parameter denotes field elements. Thus a sum over \(0,\ldots,len-1\) always ranges over decoded message blocks, even though bytes per block change with \(\pi\).

Helper functions are internal. They may take DSL values as parameters and return indices or field elements, but are not exposed through the stable C hash API.

\section{Expressions and Operations}
\label{sec:dsl-expressions}

Index expressions contain ordinary integer operations and describe loop bounds and buffer positions.

Buffer loads decode byte chunks into field elements. The load index is zero-based and counts field elements, not bytes. If \(c=\lfloor \pi/8\rfloor\), then \(message[i]\) reads bytes beginning at offset \(ic\). The exact load code is target-dependent.

The expression language contains unary negation, binary arithmetic operators \(+\), \(-\), \(*\), \(/\), \(\%\), and comparisons. Arithmetic operands must have the same DSL type. Index operands use integer arithmetic; field-element operands use arithmetic modulo the selected prime. Comparisons return index values used as conditions.

Exponentiation is restricted to \(0\), \(1\), and \(2\). Field-element \(x^2\) lowers to a dedicated square.

\section{Reductions}
\label{sec:dsl-reductions}

The DSL contains three reduction forms, summarized in \Cref{tab:reduction-forms}.

\begin{table}[t]
    \centering
    \small
    \setlength{\tabcolsep}{3pt}
    \caption{Reduction forms supported by the DSL.}
    \label{tab:reduction-forms}
    \begin{tabular}{p{0.2\textwidth}p{0.35\textwidth}p{0.3\textwidth}}
        \toprule
        Construct & Accumulator update & Vectorization \\
        \midrule
        \texttt{sum} &
        \(acc + e_i\) &
        yes \\
        \texttt{horner} &
        \(r \cdot (acc + e_i)\) &
        yes, via branches \\
        \texttt{foldl} &
        \(f(i, acc)\) &
        scalar only \\
        \bottomrule
    \end{tabular}
\end{table}

\paragraph{Sums.}
A sum has an index variable, a range, a step, and a body expression. It denotes the finite sum over the index domain. Several summands can often be evaluated independently and accumulated in parallel.

\paragraph{Horner reductions.}
A Horner reduction evaluates a polynomial recurrence. It takes a range, a body expression, and a multiplier \(r\), evaluated once before the loop:
\[
    acc \leftarrow r \cdot (acc + e_i),
\]
where \(e_i\) is the body expression at index \(i\). This captures Poly1305-style evaluation and branch-parallel vectorization~\cite{bernstein2007pema,degabriele2024sok}.

\paragraph{Left folds.}
A left fold introduces an accumulator name and initial value. For each index, the body computes the next accumulator:
\[
    acc_0 \leftarrow init,\qquad acc_{j+1} \leftarrow f(i_j, acc_j).
\]
Left folds are scalar and serve loop-carried designs such as iterative HKM.

\section{Conditionals}
\label{sec:dsl-conditionals}

The language distinguishes constant-time and non-constant-time conditionals. The regular \texttt{if} computes both branches and selects the result. The \texttt{nctif} construct lowers to an ordinary branch and may omit an else branch.

Branching on public information, such as message length parity, is usually acceptable; branching on secret data is not. The DSL makes this distinction explicit.

\section{Running example}
As a running example, consider
\[
    S(r, M,n)=\sum_{i=0}^{n-1} M_i^2,
\]
where each \(M_i\) is decoded from a message buffer. \Cref{lst:squaresum-dsl} shows the corresponding DSL module.

\begin{lstlisting}[style=thesisBNF,caption={DSL expression for the running square-sum example.},label={lst:squaresum-dsl}]
hash squaresum(key buffer, message buffer, len index) {
    sum[i,0:len:1]{message[i] ** 2}
}
\end{lstlisting}
